\documentclass[11pt,a4paper]{article}

%% ---- Packages ---------------------------------------------------------------
\usepackage[T1]{fontenc}
\usepackage[utf8]{inputenc}
\usepackage{lmodern}
\usepackage{microtype}
\usepackage{geometry}
\usepackage{amsmath,amssymb,amsthm}
\usepackage{bm}           % bold math
\usepackage{booktabs}
\usepackage{array}
\usepackage{multirow}
\usepackage{graphicx}
\usepackage{xcolor}
\usepackage{enumitem}
\usepackage{caption}
\usepackage{subcaption}
\usepackage[hidelinks]{hyperref}
\usepackage{cleveref}
\usepackage{fancyhdr}
\usepackage{titlesec}
\usepackage{parskip}
\usepackage{tcolorbox}
\tcbuselibrary{skins,breakable}

\geometry{
  top=2.5cm, bottom=2.5cm,
  left=2.5cm, right=2.5cm
}

%% ---- Styling ----------------------------------------------------------------
\definecolor{uwbblue}{RGB}{0,82,155}
\definecolor{uwbgray}{RGB}{90,90,90}
\definecolor{codegray}{RGB}{245,245,245}
\definecolor{warnbg}{RGB}{255,250,230}
\definecolor{warnborder}{RGB}{200,150,0}

\titleformat{\section}{\large\bfseries\color{uwbblue}}{{\thesection}}{1em}{}[\color{uwbblue}\titlerule]
\titleformat{\subsection}{\normalsize\bfseries\color{uwbblue}}{{\thesubsection}}{1em}{}
\titleformat{\subsubsection}{\normalsize\bfseries\color{uwbgray}}{{\thesubsubsection}}{1em}{}

\pagestyle{fancy}
\fancyhf{}
\fancyhead[L]{\small\color{uwbgray}UWB RTLS — Design Document}
\fancyhead[R]{\small\color{uwbgray}\leftmark}
\fancyfoot[C]{\small\thepage}
\renewcommand{\headrulewidth}{0.4pt}

\newtcolorbox{keybox}[1][]{
  enhanced, breakable,
  colback=codegray, colframe=uwbblue!60,
  boxrule=0.5pt, arc=3pt,
  title=#1, fonttitle=\bfseries\small\color{uwbblue},
  attach boxed title to top left={yshift=-2mm, xshift=4mm},
  boxed title style={colback=white, colframe=uwbblue!60}
}

\newtcolorbox{notebox}{
  enhanced, breakable,
  colback=warnbg, colframe=warnborder,
  boxrule=0.5pt, arc=3pt, left=4pt, right=4pt
}

%% ---- Math macros ------------------------------------------------------------
\newcommand{\vect}[1]{\bm{#1}}
\newcommand{\mat}[1]{\mathbf{#1}}
\newcommand{\norm}[1]{\left\|#1\right\|}
\newcommand{\abs}[1]{\left|#1\right|}
\newcommand{\Trans}{^{\top}}
\newcommand{\R}{\mathbb{R}}
\newcommand{\eye}[1]{\mat{I}_{#1}}
\DeclareMathOperator{\diag}{diag}
\DeclareMathOperator{\tr}{tr}
\DeclareMathOperator{\atan}{atan2}

%% ---- Document ---------------------------------------------------------------
\begin{document}

\begin{titlepage}
  \centering
  \vspace*{2cm}
  {\color{uwbblue}\rule{\linewidth}{2pt}}\\[0.5cm]
  {\huge\bfseries Ultra-Wideband Indoor Positioning System\\[0.3cm]
  \large Design Document: Motivation, Objectives, Methodology \& Mathematics}\\[0.4cm]
  {\color{uwbblue}\rule{\linewidth}{2pt}}\\[1.5cm]

  {\large\itshape Target Platform:} \textbf{Makerfabs ESP32 UWB Pro (DW1000)}\\[0.3cm]
  {\large\itshape Sensing:} \textbf{UWB Ranging + BNO085 IMU}\\[0.3cm]
  {\large\itshape Output:} \textbf{XYZ Position + Roll/Pitch/Yaw Orientation}\\[1.5cm]

  \begin{tabular}{ll}
    \textbf{Version}  & 2.0 \\
    \textbf{Date}     & June 2026 \\
    \textbf{Status}   & Active Development \\
  \end{tabular}
  \vfill
\end{titlepage}

\tableofcontents
\newpage

%% =============================================================================
\section{Introduction and Motivation}
%% =============================================================================

\subsection{Problem Statement}

Accurate indoor positioning is a prerequisite for autonomous robotics, asset tracking, human-motion capture, augmented reality, and industrial safety systems. The Global Navigation Satellite System (GNSS) fails indoors because satellite signals are attenuated by building materials and the reflected multipath environment renders pseudorange measurements unreliable. Optical systems (cameras, LiDAR) achieve high accuracy but require significant infrastructure, computational resources, and are sensitive to lighting conditions and occlusion. Inertial systems drift unboundedly without external correction.

The result is a gap: no low-cost, infrastructure-light technology currently provides real-time 3D position \emph{and} 3D orientation (six degrees of freedom, 6-DOF) indoors at centimetre-level accuracy with off-the-shelf microcontroller hardware.

\subsection{Ultra-Wideband as a Candidate Technology}

Ultra-wideband (UWB) radio occupies a spectrum of at least 500\,MHz bandwidth (typically 3.1--10.6\,GHz for IEEE 802.15.4a devices). The key physical properties that make UWB suitable for ranging are:

\begin{itemize}[leftmargin=2em]
  \item \textbf{Time resolution.} The large bandwidth produces a narrow autocorrelation peak. The Decawave DW1000 timestamps signal arrivals at 15.65\,ps resolution (64\,GHz system clock), corresponding to a theoretical distance resolution of $15.65 \times 10^{-12} \times 3 \times 10^8 \approx 4.7$\,mm per clock tick.
  \item \textbf{Multipath separation.} The short pulse duration ($\sim$2\,ns) allows the receiver to distinguish the direct-path component from reflections that arrive within a few nanoseconds. Narrowband systems (Bluetooth, Wi-Fi, Zigbee) cannot make this separation.
  \item \textbf{Material penetration.} UWB signals propagate through common building materials (drywall, wood, glass), enabling ranging through soft obstructions at the cost of a positive distance bias.
  \item \textbf{Low power.} The FCC Part 15 power limit ($-41.3$\,dBm/MHz) combined with the wide bandwidth results in pulse energy well below noise floor for narrowband receivers, making UWB inherently covert and interference-resistant.
\end{itemize}

The Decawave DW1000 chip, used in the Makerfabs ESP32 UWB Pro module, implements IEEE 802.15.4-2011 UWB and achieves a specified ranging accuracy of $\pm$10\,cm in line-of-sight (LOS) conditions.

\subsection{Limitations of Existing Open-Source Libraries}

Five open-source Arduino libraries implement DW1000 positioning. All share structural limitations that prevent scaling beyond four anchors or a single mobile tag:

\begin{table}[h!]
\centering
\caption{Open-source DW1000 library comparison}
\begin{tabular}{lccccc}
\toprule
\textbf{Feature} & \textbf{thotro} & \textbf{Makerfabs} & \textbf{jremington} & \textbf{pizzo00} & \textbf{dw1000-ng} \\
\midrule
Anchor ceiling         & 4  & 4  & 7* & 6  & Unlimited \\
RF collisions          & Yes & Yes & Yes & Yes & No \\
Update rate (Hz)       & 0.6 & 0.6 & 0.6 & 1 & 2 \\
Multi-tag              & No & No & No & Partial & No \\
NLOS detection         & No & No & No & No & No \\
3D positioning         & No & No & Fragile & No & No \\
IMU fusion             & No & No & No & No & No \\
Maintained             & No & Partial & Yes & Yes & No \\
\bottomrule
\end{tabular}
\end{table}

The four-anchor ceiling in \texttt{thotro} and its derivatives arises from the RANGE frame encoding: each device occupies 17\,bytes (2-byte address + three 5-byte absolute timestamps) within a 90-byte frame buffer. With the 9-byte MAC header, only $\lfloor(90-9-2)/17\rfloor = 4$ devices fit before buffer overflow. The \texttt{pizzo00} fork reduces this to 12\,bytes/device using differential timestamps, reaching 6 anchors, but the broadcast POLL mechanism still causes RF collision between anchor replies.

Critically, \emph{no existing library} uses the DW1000's first-path power registers for non-line-of-sight (NLOS) detection, despite the chip making this data available on every received frame.

\subsection{Proposed Contribution}

This project designs and implements a scalable UWB real-time locating system (RTLS) with the following distinguishing characteristics:

\begin{enumerate}[leftmargin=2em]
  \item \textbf{Unlimited anchor scaling} via individual addressed DS-TWR exchanges (no shared frame buffer, no broadcast collisions).
  \item \textbf{NLOS detection} using the DW1000 first-path power registers (FP\_AMPL1/2/3) — the first open-source implementation of this capability.
  \item \textbf{Physics-aware state estimation} using a Moving Horizon Estimator with Huber loss, physical room-boundary constraints, and speed limits.
  \item \textbf{Full 6-DOF output} (XYZ position + RPY orientation) via tight fusion of UWB ranging and BNO085 IMU data.
  \item \textbf{Anchor self-survey} via anchor-to-anchor ranging and Multidimensional Scaling, eliminating the need for manual coordinate measurement.

\end{enumerate}

%% =============================================================================
\section{System Overview}
%% =============================================================================

\subsection{Hardware}

\begin{keybox}[Hardware Specification]
\begin{tabular}{ll}
  \textbf{UWB module} & Makerfabs ESP32 UWB Pro with Display \\
  \textbf{UWB chip}   & Decawave DW1000 (IEEE 802.15.4-2011 UWB) \\
  \textbf{MCU}        & ESP32 (dual-core Xtensa LX6, 240\,MHz, Wi-Fi/BT) \\
  \textbf{RF front-end} & PA/LNA amplifier ($\sim$200\,m outdoor range) \\
  \textbf{Display}    & 1.3\,inch SSD1306 OLED (I\textsuperscript{2}C, pins SDA=4, SCL=5) \\
  \textbf{IMU (planned)} & Bosch BNO085 via SPI \\
  \textbf{Configuration} & 3--4 anchors (fixed) + 1 tag (mobile) \\
  \textbf{Transport}  & Wi-Fi UDP to MATLAB host (192.168.x.x:5005) \\
\end{tabular}
\end{keybox}

\subsection{Software Architecture}

The system is split into two layers:

\textbf{Firmware (Arduino/ESP32):} The tag performs round-robin DS-TWR ranging to each anchor, collects raw distances with quality metrics, and streams measurement packets over UDP. Anchors are purely passive responders. The tag never solves its own position.

\textbf{Host (MATLAB):} All position computation, filtering, visualisation, and logging runs on the host PC. This separation means adding anchors requires only editing a JSON configuration file — no firmware change. The firmware packet format is versioned and IMU-ready.

The processing pipeline on the host is:
\[
\text{Raw ranges} \xrightarrow{\text{median filter}} \text{Filtered ranges} \xrightarrow{\text{NLOS weighting}} \text{Weighted ranges} \xrightarrow{\text{LM multilat.}} \text{Position} \xrightarrow{\text{EKF/MHE}} \hat{\vect{x}} \xrightarrow{\text{EMA}} \text{Output}
\]

%% =============================================================================
\section{Objectives}
%% =============================================================================

\begin{enumerate}[leftmargin=2em, label=\textbf{O\arabic*.}]
  \item Achieve XY positioning accuracy of $\leq 3$\,cm (1-$\sigma$) in LOS and $\leq 8$\,cm in soft NLOS conditions using 3--4 anchors.
  \item Achieve Z-axis positioning accuracy of $\leq 5$\,cm with 4 anchors deployed at different heights.
  \item Deliver continuous RPY orientation output at $\geq 50$\,Hz from BNO085 IMU fusion.
  \item Maintain position tracking continuity during brief anchor occlusions ($\leq 3$\,s) using IMU dead-reckoning.
  \item Support anchor self-survey: determine anchor positions to $\leq 2$\,cm without external measurement.
  \item Demonstrate update rate $\geq 3$\,Hz for 3 anchors and $\geq 8$\,Hz in fast-mode operation.
  \item Provide real-time NLOS detection per anchor using DW1000 channel impulse response diagnostics.
  \item Remain compatible with low-cost hardware: total cost $<$ \$200 for a 4-anchor + 1-tag system.
\end{enumerate}

%% =============================================================================
\section{Mathematical Foundations}
%% =============================================================================

\subsection{Double-Sided Two-Way Ranging (DS-TWR)}

\subsubsection{Motivation}

Simple one-way ranging computes distance from the time of flight (ToF) of a single transmitted pulse. This requires synchronised clocks between transmitter and receiver. Crystal oscillators in consumer devices have frequency offsets of $\pm20$\,ppm; at 7\,m range, this alone introduces a ranging error of $\pm 0.7$\,mm per millisecond of response latency — growing to centimetres with longer delays.

Double-Sided Two-Way Ranging eliminates the clock-offset term analytically through a four-message exchange that gives the receiver two independent estimates of the one-way ToF.

\subsubsection{Message Exchange}

The DS-TWR protocol involves four timestamped events across tag (T) and anchor (A):

\begin{center}
\begin{tabular}{clll}
\toprule
\textbf{Step} & \textbf{Direction} & \textbf{Frame} & \textbf{Timestamps captured} \\
\midrule
1 & T $\to$ A & \textsc{poll}      & $t_1$ (T tx), $t_2$ (A rx) \\
2 & A $\to$ T & \textsc{poll-ack}  & $t_3$ (A tx), $t_4$ (T rx) \\
3 & T $\to$ A & \textsc{range}     & $t_5$ (T tx); payload: $t_1,t_4,t_5$ \\
4 & A $\to$ T & \textsc{range-report} & $t_6$ (A rx); anchor computes \& replies $d$ \\
\bottomrule
\end{tabular}
\end{center}

\subsubsection{ToF Formula}

Define the following intervals (all in DW1000 clock ticks, 15.65\,ps each):
\begin{align}
R_a &= t_4 - t_1 & &\text{(tag: total round-trip, poll to ack-received)} \\
D_a &= t_3 - t_2 & &\text{(anchor: processing delay, poll-rx to ack-tx)} \\
R_b &= t_6 - t_3 & &\text{(anchor: total round-trip, ack to range-received)} \\
D_b &= t_5 - t_4 & &\text{(tag: processing delay, ack-rx to range-tx)}
\end{align}

The asymmetric DS-TWR formula (Decawave APS013) gives the time-of-flight:
\begin{equation}
\boxed{
\tau_\text{ToF} = \frac{R_a R_b - D_a D_b}{R_a + R_b + D_a + D_b}
}
\label{eq:dstwr}
\end{equation}

\textbf{Clock-offset cancellation.} Let $\varepsilon_T, \varepsilon_A$ be the fractional frequency offsets of the tag and anchor crystals respectively. Including these offsets, the true one-way ToF $\tau$ satisfies:
\[
R_a = (2\tau + D_a)(1 + \varepsilon_T), \quad R_b = (2\tau + D_b)(1 + \varepsilon_A)
\]
Substituting into \cref{eq:dstwr} and expanding to first order in $\varepsilon$, the frequency-offset terms cancel exactly, leaving:
\[
\tau_\text{ToF} \approx \tau + O(\varepsilon^2)
\]
Residual error scales as $\varepsilon^2 \approx (20\,\text{ppm})^2 = 4 \times 10^{-10}$, contributing $<0.1$\,mm of ranging error for exchange times $\leq 10$\,ms.

The measured distance is:
\begin{equation}
\hat{d} = c \cdot \tau_\text{ToF} + b_\text{ant}
\label{eq:dist}
\end{equation}
where $c = 2.998 \times 10^8$\,m/s and $b_\text{ant}$ is the per-board antenna delay calibration offset.

\subsubsection{Antenna Delay Calibration}

The DW1000 applies a programmable antenna delay to the RX timestamp register. If the true distance is $d^\star$ and the raw measured distance is $\hat{d}(\Delta)$ for delay value $\Delta$, calibration solves:
\[
\Delta^* = \arg\min_\Delta \left(\mathbb{E}\left[\hat{d}(\Delta)\right] - d^\star\right)^2
\]
via binary search. The relationship is monotone: increasing $\Delta$ shifts the RX timestamp later, reducing the computed ToF and thus the reported distance. A binary search over $\Delta \in [15800, 16900]$ with 200 averaged measurements per step converges to the optimal delay within $\pm 3$ ticks ($\approx \pm 1.4$\,mm) in 14 iterations.

\subsection{NLOS Detection via Channel Impulse Response}

\subsubsection{Physical Basis}

The DW1000 correlates the received signal against a known preamble sequence to produce a Channel Impulse Response (CIR) estimate $h(\tau)$. In LOS conditions the CIR has a single dominant peak at the direct-path delay $\tau_0$. In NLOS conditions (blocked or obstructed direct path), reflected components at delays $\tau_0 + \delta_k$ contribute additional energy to the CIR:
\[
h(\tau) = \alpha_0 \delta(\tau - \tau_0) + \sum_{k=1}^{K} \alpha_k \delta(\tau - \tau_0 - \delta_k)
\]

The total received power integrates over all components; the first-path power reflects only the direct component.

\subsubsection{DW1000 Diagnostic Registers}

After each receive, the DW1000 populates:
\begin{itemize}[leftmargin=2em]
  \item \textbf{FP\_AMPL1, FP\_AMPL2, FP\_AMPL3} (register 0x12, offsets 0x07--0x0C): amplitude of the three samples around the first-path leading edge.
  \item \textbf{RXPACC} (register 0x10, offset 0x2C): accumulated path count (preamble accumulator length).
  \item \textbf{CIR\_PWR} (register 0x12, offset 0x00): total channel power from the accumulator.
\end{itemize}

\subsubsection{NLOS Score Computation}

The first-path power (Decawave APS006 formula) is:
\begin{equation}
P_\text{FP} = 10 \log_{10}\!\left(\frac{A_1^2 + A_2^2 + A_3^2}{N_\text{acc}^2}\right) - A_\text{const}
\label{eq:fppower}
\end{equation}
where $A_1, A_2, A_3$ are the FP\_AMPL register values, $N_\text{acc}$ is RXPACC, and $A_\text{const} = 121.74$\,dBm for 16\,MHz PRF or $113.77$\,dBm for 64\,MHz PRF.

The total received power is computed from CIR\_PWR by the existing driver function \texttt{DW1000.getReceivePower()}.

The NLOS score is the power differential:
\begin{equation}
\boxed{
\Delta P_\text{NLOS} = P_\text{RX} - P_\text{FP}
}
\label{eq:nlos}
\end{equation}

\textbf{Interpretation:}
\begin{itemize}[leftmargin=2em]
  \item $\Delta P_\text{NLOS} < 3$\,dB: LOS. Direct path dominates. Full measurement weight.
  \item $3 \leq \Delta P_\text{NLOS} < 6$\,dB: Soft NLOS. Partial obstruction or light multipath. Reduce weight.
  \item $\Delta P_\text{NLOS} \geq 6$\,dB: Hard NLOS. Reflected path dominant. Strong downweighting or rejection.
\end{itemize}

\subsubsection{Range Bias under NLOS}

NLOS ranging error is always a \emph{positive} bias: reflected and through-wall paths are geometrically longer than the direct path. Through-wall propagation in typical building materials also slows the signal by a factor $\sqrt{\varepsilon_r}$ where $\varepsilon_r$ is the relative permittivity (concrete: $\varepsilon_r \approx 6$--$9$; drywall: $\varepsilon_r \approx 2$--$4$). The empirical range bias $b_\text{NLOS}$ correlates with $\Delta P_\text{NLOS}$:
\begin{equation}
b_\text{NLOS} \approx k_\text{NLOS} \cdot \max(0,\; \Delta P_\text{NLOS} - 3\,\text{dB})
\label{eq:biasnlos}
\end{equation}
with $k_\text{NLOS} \approx 0.02$--$0.05$\,m/dB (environment-dependent, tunable). This correction is applied before multilateration.

\subsection{Weighted Multilateration}

\subsubsection{Measurement Model}

Given $N$ anchors at known positions $\vect{a}_i \in \R^d$ ($d = 2$ or $3$) and measured ranges $\hat{d}_i$ (after median filtering and bias correction), the tag position $\vect{x} \in \R^d$ satisfies:
\begin{equation}
\hat{d}_i = \norm{\vect{x} - \vect{a}_i} + e_i, \quad e_i \sim \mathcal{N}(0, \sigma_i^2)
\label{eq:measurement}
\end{equation}
The noise standard deviation $\sigma_i$ is range-quality-dependent:
\begin{equation}
\sigma_i = \sigma_\text{base} \cdot 10^{\Delta P_{\text{NLOS},i} / 20}
\label{eq:sigma}
\end{equation}
so that a 10\,dB NLOS score increases the assumed noise by a factor of $\sqrt{10} \approx 3.16$ and reduces the weight by a factor of 10.

\subsubsection{Weighted Least-Squares Formulation}

The weighted multilateration objective is:
\begin{equation}
\hat{\vect{x}} = \arg\min_{\vect{x} \in \R^d} \sum_{i=1}^{N} w_i \left(\norm{\vect{x} - \vect{a}_i} - \hat{d}_i\right)^2
\label{eq:wls}
\end{equation}
with weights $w_i = 1/\sigma_i^2$.

This is a nonlinear least-squares problem solved by Levenberg-Marquardt (LM). Define the residual vector $\vect{r}(\vect{x}) \in \R^N$ with $r_i = \norm{\vect{x} - \vect{a}_i} - \hat{d}_i$, the Jacobian $\mat{J} \in \R^{N \times d}$ with rows:
\begin{equation}
\mat{J}_i = \frac{\partial r_i}{\partial \vect{x}} = \frac{(\vect{x} - \vect{a}_i)\Trans}{\norm{\vect{x} - \vect{a}_i}}
\label{eq:jacobian}
\end{equation}
and the weight matrix $\mat{W} = \diag(w_1, \ldots, w_N)$.

The LM update step is:
\begin{equation}
\vect{x}_{k+1} = \vect{x}_k - \left(\mat{J}\Trans \mat{W} \mat{J} + \lambda \diag(\mat{J}\Trans \mat{W} \mat{J})\right)^{-1} \mat{J}\Trans \mat{W} \vect{r}(\vect{x}_k)
\label{eq:lm}
\end{equation}
The damping parameter $\lambda$ is increased when a step increases cost and decreased otherwise, interpolating between gradient descent (large $\lambda$) and Gauss-Newton (small $\lambda$).

\subsubsection{Robust Outlier Gating}

After a first LM solve, outlier anchors are identified and gated. The residuals $\{r_i\}$ are assessed against a scaled median absolute deviation (MAD):
\begin{equation}
\hat{\sigma}_\text{robust} = 1.4826 \cdot \text{median}\left(\abs{r_i - \text{median}(r_i)}\right)
\end{equation}
Anchors with $\abs{r_i - \text{median}(r_i)} > K \cdot \hat{\sigma}_\text{robust}$ (default $K = 2.5$) are excluded and the solve is repeated with the remaining anchors. The factor 1.4826 makes the MAD a consistent estimator of $\sigma$ under a Gaussian assumption.

\subsubsection{Covariance Approximation}

The position covariance is approximated from the Hessian of the gated objective at the solution:
\begin{equation}
\mat{P}_\text{pos} \approx \sigma_\text{base}^2 \left(\mat{J}\Trans \mat{W} \mat{J}\right)^{-1}
\label{eq:cov}
\end{equation}
This is passed to the Kalman filter as the measurement covariance $\mat{R}_k$.

\subsubsection{GPS-Style Simultaneous Bias Estimation}

When $N \geq d+2$ anchors are available, the augmented state $\tilde{\vect{x}} = [\vect{x}\Trans, \delta]\Trans \in \R^{d+1}$ is solved, where $\delta$ is a common range bias (absorbing tag delay miscalibration and average multipath):
\begin{equation}
\hat{d}_i = \norm{\vect{x} - \vect{a}_i} + \delta + e_i
\end{equation}

The augmented Jacobian appends a column of ones:
\begin{equation}
\tilde{\mat{J}}_i = \left[\mat{J}_i,\; 1\right]
\label{eq:jacaug}
\end{equation}

This is directly analogous to GPS pseudorange processing, where the receiver clock offset appears as a common additive term on all satellite measurements.

\begin{notebox}
\textbf{Geometric note.} Solving for $\delta$ with $N = d+1$ anchors (minimum for standard multilateration) is equivalent to Time Difference of Arrival (TDOA) positioning, since $\delta$ cancels in all pairwise range differences. TDOA is geometrically weaker than TOA for symmetric anchor layouts — the problem becomes ill-conditioned near the anchor centroid. This mode is therefore only enabled when $N \geq d+2$, providing at least one degree of overdetermination.
\end{notebox}

\subsection{State Estimation}

\subsubsection{Extended Kalman Filter}

The EKF models the tag as a constant-acceleration (CA) particle. The state vector is:
\begin{equation}
\vect{x} = \begin{bmatrix} \vect{p} \\ \vect{v} \\ \vect{a} \\ \delta \end{bmatrix} \in \R^{3d+1}
\end{equation}
where $\vect{p} \in \R^d$ is position, $\vect{v} \in \R^d$ is velocity, $\vect{a} \in \R^d$ is acceleration, and $\delta \in \R$ is the slowly-varying range bias.

The discrete-time transition model with time step $\Delta t$ is:
\begin{equation}
\vect{x}_{k+1} = \mat{F} \vect{x}_k + \vect{w}_k, \quad \vect{w}_k \sim \mathcal{N}(\vect{0}, \mat{Q})
\end{equation}
\begin{equation}
\mat{F} = \begin{bmatrix}
\mat{I}_d & \Delta t \mat{I}_d & \tfrac{\Delta t^2}{2}\mat{I}_d & \vect{0} \\
\mat{0}   & \mat{I}_d         & \Delta t \mat{I}_d              & \vect{0} \\
\mat{0}   & \mat{0}           & \mat{I}_d                       & \vect{0} \\
0         & 0                 & 0                               & 1
\end{bmatrix}
\end{equation}

The process noise covariance for the CA model is (Singer model):
\begin{equation}
\mat{Q} = q_j \begin{bmatrix}
\tfrac{\Delta t^5}{20}\mat{I} & \tfrac{\Delta t^4}{8}\mat{I} & \tfrac{\Delta t^3}{6}\mat{I} & \vect{0} \\
\tfrac{\Delta t^4}{8}\mat{I}  & \tfrac{\Delta t^3}{3}\mat{I} & \tfrac{\Delta t^2}{2}\mat{I} & \vect{0} \\
\tfrac{\Delta t^3}{6}\mat{I}  & \tfrac{\Delta t^2}{2}\mat{I} & \Delta t \mat{I}              & \vect{0} \\
\vect{0} & \vect{0} & \vect{0} & q_\delta \Delta t
\end{bmatrix}
\end{equation}
where $q_j$ is the jerk power spectral density (tuning parameter) and $q_\delta$ governs the rate of bias drift.

The measurement model for anchor $i$ is (tight coupling):
\begin{equation}
h_i(\vect{x}) = \norm{\vect{p} - \vect{a}_i} + \delta
\label{eq:measmodel}
\end{equation}
with Jacobian row:
\begin{equation}
\mat{H}_i = \left[\frac{(\vect{p} - \vect{a}_i)\Trans}{\norm{\vect{p} - \vect{a}_i}},\; \vect{0}_{1 \times 2d},\; 1\right]
\end{equation}

The standard Kalman prediction and update equations are:
\begin{align}
\text{Predict:}\quad & \hat{\vect{x}}_{k|k-1} = \mat{F}\hat{\vect{x}}_{k-1|k-1} \\
                     & \mat{P}_{k|k-1} = \mat{F}\mat{P}_{k-1|k-1}\mat{F}\Trans + \mat{Q} \\
\text{Update:}\quad  & \mat{S}_k = \mat{H}_k \mat{P}_{k|k-1} \mat{H}_k\Trans + \mat{R}_k \\
                     & \mat{K}_k = \mat{P}_{k|k-1} \mat{H}_k\Trans \mat{S}_k^{-1} \\
                     & \hat{\vect{x}}_{k|k} = \hat{\vect{x}}_{k|k-1} + \mat{K}_k \left(\vect{z}_k - h(\hat{\vect{x}}_{k|k-1})\right) \\
                     & \mat{P}_{k|k} = (\mat{I} - \mat{K}_k \mat{H}_k)\mat{P}_{k|k-1}
\end{align}

\subsubsection{Moving Horizon Estimator}

The MHE solves a constrained nonlinear optimisation over a sliding window of $N$ past measurement epochs. At time $k$, the window spans $\{k-N+1, \ldots, k\}$.

\textbf{Objective function:}
\begin{equation}
\min_{\{\vect{x}_t\}_{t=k-N+1}^{k}} \underbrace{\norm{\vect{x}_{k-N+1} - \bar{\vect{x}}_{k-N+1}}_{\mat{P}_\text{arr}^{-1}}^2}_{\text{arrival cost}} + \underbrace{\sum_{t=k-N+2}^{k} \norm{\vect{x}_t - \mat{F}\vect{x}_{t-1}}_{\mat{Q}^{-1}}^2}_{\text{dynamics residuals}} + \underbrace{\sum_{t=k-N+1}^{k} \sum_{i=1}^{N_a} \rho_\delta\!\left(\norm{\vect{p}_t - \vect{a}_i} - \hat{d}_{t,i}\right)}_{\text{range residuals}}
\label{eq:mhe}
\end{equation}

where $\norm{\vect{v}}_{\mat{M}}^2 = \vect{v}\Trans \mat{M} \vect{v}$, $\bar{\vect{x}}_{k-N+1}$ and $\mat{P}_\text{arr}$ are the prior mean and covariance from the previous window solution (arrival cost captures past information compactly), and $\rho_\delta$ is the Huber loss function:
\begin{equation}
\rho_\delta(e) = \begin{cases}
\dfrac{1}{2}e^2 & \text{if } \abs{e} \leq \delta \\[6pt]
\delta\!\left(\abs{e} - \dfrac{\delta}{2}\right) & \text{if } \abs{e} > \delta
\end{cases}
\label{eq:huber}
\end{equation}
with threshold $\delta = 0.15$\,m (one standard deviation of typical ranging noise). Large residuals from NLOS anchors grow linearly rather than quadratically, reducing their influence without complete rejection.

\textbf{Constraints:}
\begin{align}
\vect{x}_\text{min} &\leq \vect{p}_t \leq \vect{x}_\text{max} & &\text{(room boundary, hard)} \label{eq:bound}\\
\norm{\vect{v}_t}_2 &\leq v_\text{max} & &\text{(speed limit, hard)} \label{eq:speed}
\end{align}

The boundary constraint prevents the estimator from producing physically impossible positions during anchor occlusion. The speed limit (e.g.\ $v_\text{max} = 2$\,m/s for walking) eliminates large transient jumps.

\textbf{Computational cost.} With $N = 6$, $d = 3$, and $N_a = 4$ anchors, the optimisation has $N \times (3d+1) = 60$ decision variables and is solved by \texttt{fmincon} (interior-point method) in approximately 2--5\,ms per step, well within the $\sim$333\,ms UWB sweep period.

\textbf{Comparison with EKF:}
\begin{table}[h!]
\centering
\caption{EKF vs MHE comparison for this application}
\begin{tabular}{lll}
\toprule
\textbf{Property} & \textbf{EKF} & \textbf{MHE} \\
\midrule
NLOS handling & Pre-filtering required & Huber loss absorbs outliers \\
Physical constraints & None & Room boundary, speed limit (hard) \\
Nonlinearity & Linearised via Jacobian & Solved exactly \\
Lag & Zero (one sample) & $N/2$ samples ($\approx 1$\,s at $N=6$) \\
Computation per step & $O(d^2)$, trivial & 2--5\,ms (fmincon) \\
Sensitivity to init & Moderate & Low (window smooths errors) \\
\bottomrule
\end{tabular}
\end{table}

\subsection{IMU Fusion and Orientation Estimation}

\subsubsection{Quaternion Representation}

Orientation is represented by a unit quaternion $\vect{q} = [q_w, q_x, q_y, q_z]\Trans \in \R^4$ satisfying $\norm{\vect{q}} = 1$. The rotation matrix from body frame to world frame is:
\begin{equation}
\mat{R}(\vect{q}) = \begin{bmatrix}
1 - 2(q_y^2 + q_z^2) & 2(q_x q_y - q_w q_z) & 2(q_x q_z + q_w q_y) \\
2(q_x q_y + q_w q_z) & 1 - 2(q_x^2 + q_z^2) & 2(q_y q_z - q_w q_x) \\
2(q_x q_z - q_w q_y) & 2(q_y q_z + q_w q_x) & 1 - 2(q_x^2 + q_y^2)
\end{bmatrix}
\end{equation}

\subsubsection{Quaternion Kinematics}

The continuous-time quaternion differential equation driven by body angular rate $\vect{\omega}^b = [\omega_x, \omega_y, \omega_z]\Trans$ is:
\begin{equation}
\dot{\vect{q}} = \frac{1}{2} \vect{q} \otimes \begin{bmatrix}0 \\ \vect{\omega}^b\end{bmatrix}
\label{eq:qdot}
\end{equation}
where $\otimes$ denotes the Hamilton product. The discrete propagation over $\Delta t$ is:
\begin{equation}
\vect{q}_{k+1} = \vect{q}_k \otimes \exp\!\left(\frac{\Delta t}{2}\begin{bmatrix}0 \\ \vect{\omega}^b_k\end{bmatrix}\right) \approx \vect{q}_k \otimes \begin{bmatrix}\cos(\frac{\Delta t}{2}\norm{\vect{\omega}^b_k}) \\ \frac{\vect{\omega}^b_k}{\norm{\vect{\omega}^b_k}}\sin(\frac{\Delta t}{2}\norm{\vect{\omega}^b_k})\end{bmatrix}
\label{eq:qprop}
\end{equation}

\subsubsection{IMU-Driven Position Propagation}

The BNO085 outputs linear acceleration $\vect{a}^b$ (gravity removed) in the body frame. The world-frame acceleration is $\vect{a}^w = \mat{R}(\vect{q}) \vect{a}^b$. Position and velocity propagate as:
\begin{align}
\vect{p}_{k+1} &= \vect{p}_k + \vect{v}_k \Delta t + \tfrac{1}{2} \vect{a}^w_k \Delta t^2 \\
\vect{v}_{k+1} &= \vect{v}_k + \vect{a}^w_k \Delta t
\end{align}

This propagation runs at the IMU rate (100\,Hz), providing 33 intermediate position estimates between each UWB measurement update (3\,Hz). UWB range measurements correct the accumulated drift at each available sweep.

\subsubsection{Loose vs Tight Coupling}

\textbf{Loose coupling} treats the multilaterated position as the measurement. The EKF state is updated with $\vect{z}_k = \hat{\vect{p}}_\text{multilat}$ and $\mat{R}_k = \mat{P}_\text{pos}$ from \cref{eq:cov}. This is simpler to implement but suffers from the two-stage sequential estimation error.

\textbf{Tight coupling} uses the raw UWB ranges directly as measurements via the nonlinear model of \cref{eq:measmodel}. This avoids the intermediate multilateration step, correctly propagates uncertainty, and allows the filter to use partial anchor sets (e.g. 2 of 4 anchors visible) without changing the measurement model structure.

\subsubsection{Roll, Pitch, and Yaw Extraction}

The Euler angles (roll $\phi$, pitch $\theta$, yaw $\psi$) in ZYX convention are extracted from the quaternion as:
\begin{align}
\phi   &= \text{atan2}\!\left(2(q_w q_x + q_y q_z),\; 1 - 2(q_x^2 + q_y^2)\right) \label{eq:roll}\\
\theta &= \arcsin\!\left(2(q_w q_y - q_z q_x)\right) \label{eq:pitch}\\
\psi   &= \text{atan2}\!\left(2(q_w q_z + q_x q_y),\; 1 - 2(q_y^2 + q_z^2)\right) \label{eq:yaw}
\end{align}

The BNO085 ARVR Stabilised Rotation Vector mode applies sensor fusion (accelerometer + gyroscope + magnetometer) internally at 400\,Hz and outputs a quaternion robust to magnetic disturbances. The magnetometer provides an absolute yaw reference, preventing the unbounded yaw drift that affects gyroscope-only integration.

\subsection{Anchor Self-Survey via Multidimensional Scaling}

\subsubsection{Problem Statement}

Given $N$ anchors, measure all $\binom{N}{2}$ pairwise distances $\hat{d}_{ij}$ by having each anchor temporarily act as a ranging initiator. Recover the $N$ anchor positions $\{\vect{a}_i\}$ that are consistent with these distances.

\subsubsection{Classical Multidimensional Scaling}

Construct the squared-distance matrix $\mat{D}^{(2)} \in \R^{N \times N}$ with $D^{(2)}_{ij} = \hat{d}_{ij}^2$. Apply double-centering via the centering matrix $\mat{J} = \mat{I}_N - \frac{1}{N}\vect{1}\vect{1}\Trans$:
\begin{equation}
\mat{B} = -\tfrac{1}{2} \mat{J} \mat{D}^{(2)} \mat{J}
\label{eq:mds}
\end{equation}

If the true distances were exact, $\mat{B}$ would be a rank-$d$ positive semidefinite matrix satisfying $\mat{B} = \mat{X}\mat{X}\Trans$ where $\mat{X} \in \R^{N \times d}$ contains the true centred positions. With noisy distances, $\mat{B}$ is full rank. The best rank-$d$ approximation is obtained from the eigendecomposition $\mat{B} = \mat{V}\mat{\Lambda}\mat{V}\Trans$ by taking the $d$ largest eigenvalues:
\begin{equation}
\hat{\mat{X}} = \mat{V}_d \mat{\Lambda}_d^{1/2}
\label{eq:mdssoln}
\end{equation}
where $\mat{V}_d \in \R^{N \times d}$ and $\mat{\Lambda}_d = \diag(\lambda_1, \ldots, \lambda_d)$ contain the $d$ dominant eigenpairs.

\subsubsection{Coordinate Frame Fixation}

The MDS solution is defined only up to a rigid-body transformation (rotation, translation, reflection). Fix the coordinate frame by convention:
\begin{enumerate}[leftmargin=2em]
  \item \textbf{Translation:} Subtract $\hat{\vect{x}}_1$ (row 1 of $\hat{\mat{X}}$) from all rows so anchor 1 is at the origin.
  \item \textbf{Rotation:} Rotate so that $\hat{\vect{x}}_2 - \hat{\vect{x}}_1$ aligns with the positive X axis.
  \item \textbf{Reflection:} If anchor 3 lies in the negative Y half-plane, negate all Y coordinates.
\end{enumerate}

The resulting layout is written to \texttt{config/anchors.json} and used for all subsequent localisation without modification.

\subsubsection{Noise Propagation}

Each pairwise distance is the mean of $M = 100$ ranging measurements. If individual measurement noise is $\sigma_d$, the mean has uncertainty $\sigma_d/\sqrt{M}$. For $\sigma_d \approx 5$\,cm and $M = 100$, the per-pair uncertainty is $0.5$\,cm. Classical MDS propagates this to anchor position uncertainty of approximately $\sigma_d/\sqrt{M}$ in each coordinate — consistent with the objective of $\leq 2$\,cm anchor self-survey accuracy.

%% =============================================================================
\section{Implementation Roadmap}
%% =============================================================================

The system is developed in six phases, each independently testable:

\begin{table}[h!]
\centering
\caption{Implementation phases and key items}
\begin{tabular}{cllc}
\toprule
\textbf{Phase} & \textbf{Item} & \textbf{Description} & \textbf{Difficulty} \\
\midrule
\multirow{6}{*}{1 — Foundation}
 & 1.0 & Switch to 64\,MHz PRF (accuracy mode) & Easy \\
 & 1.1 & Better calibration (200 samples, multi-distance) & Easy \\
 & 1.2 & Adaptive polling: skip dead anchors & Easy \\
 & 1.3 & Faster sweep rate (reduce reply delay) & Medium \\
 & 1.5 & Anchor self-survey + auto \texttt{anchors.json} & Medium \\
\midrule
\multirow{5}{*}{2 — Ranging}
 & 2.1 & FP\_POWER NLOS detection (firmware + wire format) & Medium \\
 & 2.2 & Weighted multilateration (NLOS score) & Medium \\
 & 2.3 & NLOS range bias correction slider & Easy \\
 & 2.4 & Per-anchor diagnostics panel & Easy \\
 & 2.5 & GPS-style $\delta$-solve (4+ anchors) & Medium \\
\midrule
\multirow{4}{*}{3 — Estimation}
 & 3.1 & Constant-acceleration EKF & Medium \\
 & 3.2 & $\delta$ as EKF state (bias estimation) & Med-Hard \\
 & 3.3 & Moving Horizon Estimator & Hard \\
 & 3.4 & Hybrid EKF + MHE & Medium \\
\midrule
\multirow{5}{*}{4 — 3D + RPY}
 & 4.1 & 4th anchor + 3D configuration & Easy \\
 & 4.2 & BNO085 SensorImu (SPI) & Medium \\
 & 4.3 & RPY output from quaternion & Easy \\
 & 4.4 & Loose coupling: IMU-aided EKF & Hard \\
 & 4.5 & Tight coupling: raw ranges + IMU in MHE & Very Hard \\
\midrule
\multirow{4}{*}{5 — Hardening}
 & 5.1 & Deep sleep on anchors & Medium \\
 & 5.2 & DW1000 driver error handling & Medium \\
 & 5.3 & Temperature compensation & Medium \\
 & 5.4 & Online self-calibration & Hard \\
\midrule
\multirow{3}{*}{6 — Scale}
 & 6.1 & Multi-tag TDMA superframe & Hard \\
 & 6.2 & Differential timestamp encoding & Hard \\
 & 6.3 & On-device 2D solve (OLED) & Medium \\
\bottomrule
\end{tabular}
\end{table}

%% =============================================================================
\section{Expected Performance}
%% =============================================================================

\begin{table}[h!]
\centering
\caption{Expected accuracy progression through implementation phases}
\begin{tabular}{clllll}
\toprule
\textbf{After Phase} & \textbf{XY accuracy} & \textbf{Z accuracy} & \textbf{RPY} & \textbf{Rate} & \textbf{Key enabler} \\
\midrule
Current & 3--10\,cm & --- & --- & 3\,Hz & Baseline \\
Phase 1 & 2--6\,cm  & --- & --- & 4--8\,Hz & Calibration + PRF \\
Phase 2 & 1--4\,cm  & --- & --- & 4--8\,Hz & NLOS detection \\
Phase 3 & 1--3\,cm  & --- & --- & 4--8\,Hz & MHE + CA-EKF \\
Phase 4 & 1--3\,cm  & 2--5\,cm & 3--5° & 100\,Hz RPY & IMU fusion \\
Phase 5 & 1--3\,cm  & 2--5\,cm & 3--5° & 100\,Hz RPY & Reliability \\
\bottomrule
\end{tabular}
\end{table}

The theoretical floor for the DW1000 chip is approximately 1--2\,cm in perfect LOS conditions (limited by 15.65\,ps timestamp resolution and residual crystal asymmetry). Practical indoor environments with multipath and occasional NLOS will see 2--5\,cm as the sustainable floor. IMU integration does not improve the UWB ranging accuracy but eliminates position lag between sweeps and provides orientation at sensor rate.

\end{document}
